% !TEX TS-program = lualatexmk
\def\zz{1}
\if1\zz\documentclass[11pt]{article}  
\else\documentclass[11pt]{scrartcl}  
\def\tt{\normalfont\ttfamily}
\fi
\usepackage{iftex}
\usepackage[margin=1.1in]{geometry}
\usepackage[parfill]{parskip}
\usepackage[p]{newpxtext}
%\usepackage{textcomp}
\usepackage[npx,supscolor=red!70!black]{superiors}
%\usepackage[uprightGreek]{newpxmath}
%\usepackage[perpage
%%,flushmargin
%]{footmisc}
\title{Footnote examples\footnote{with newpxtext}}
\author{Michael Sharpe\thanks{Thanks much.}}
\pagestyle{empty}
\begin{document}
%\traceon
\maketitle
%\traceoff
\thispagestyle{empty}
%In\footnote{Testing} this note...\textinf{12}
%
%X{\infstyle%\addfontfeatures{RawFeature=+subs}
%2}{\destyle4}
%X{\addfontfeatures{VerticalPosition=Denominator}2}

%X\textsuperscript{2}
%\end{document}

%\makeatletter
%{\sustyle8}\normalfont{\textsuperscript{8\textdagger}}In %
%{\addfontfeature{VerticalPosition=Superior}8\textdagger}\textdagger  
%{\addfontfeature{RawFeature=+smcp}AzX}
%\makeatletter
%%\show\@makefntext
%\makeatother
%\traceon\textsuperscript{\textasteriskcentered} \traceoff

The standard LaTeX document classes {\tt article}, {\tt report} and {\tt book}, the macro \verb|\maketitle| \verb|\let|s \verb|\footnote| to \verb|\thanks| and footnote marks are taken from a specified list of symbols. (In the AMS classes\footnote{amsart, amsproc, amsbook} the title footnote marks will appear only in the footnote itself at the bottom of the page.)\\

\centerline{\fbox{\begin{minipage}{4in}
Footnotes after $\backslash$maketitle use by default numeric markers beginning at 1, in some form\footnote{Best not using scaled-down lining figures.}  smaller than lining numbers. (Footnotes in minipages use lowercase letters\footnote{italic except in KOMA classes} as markers.)
\end{minipage}}
}
This footnote page snippet was prepared with {\tt lualatex} using the preamble:
\begin{verbatim}
\documentclass[11pt]{article}  
\usepackage{trace}
\usepackage[margin=1.1in]{geometry}
\usepackage[parfill]{parskip}
\usepackage[p]{newpxtext} % larger sups figures
\usepackage[npx,supscolor=red!70!black]{superiors}
\title{Footnote examples\footnote{with newpxtext}}
\author{Michael Sharpe\thanks{Thanks much.}}
\end{verbatim}

%As the {\tt defaultsups} option was not specified, footnote markers were drawn from the superior symbols and figures in the text font.

%By setting the options {\tt supscale} and {\tt supsraised}, the footnote markers were made a bit more prominent while keeping their tops about the same as without the options.
% (The superior numbers in this font have baseline at .350{\tt em} and top at .763{\tt em}. Scaling up by the factor 1.15 increases the top by 15\% of .763{\tt em} = .114{\tt em} and moves the baseline up to 1.15 \textasteriskcentered\ 350{\tt em} = 403{\tt em}.) It is not a good idea to use the {\tt supscolor} option if you are using {\tt hyperref}.
 
\end{document}
